%% kapa: exam=2 round=30 year=2019 period=1 subject=appraisal_practice subject_ko=감정평가실무 seq=4 number=4 points=10
%% source: Q-Net 감정평가사 기출문제 (cst00309) artlSeq=5205595 file=raw/1008/5205595/감정평가실무.pdf
%% conversion: kordoc 4.24.51 -> tools/gen2.py (auto); review=proofread-sonnet; flags=none
\begin{longquestion}{【문제 4】}
1년 전 임대차계약이 체결되어 있는 오피스텔에 대하여 동 임대차계약을 인수하는 조건으로 매매가 이루어졌다. 매매가는 시장의 전형적인 수익률 5.0\%를 기준으로 산출되는 임대권의 가치를 기준으로 결정되었다. 임대차계약 내용은, 계약기간 5년으로 연간 순임대료(순영업소득)는 2,200만원이고 계약기간 중 임대차조건의 변경이 없다는 내용으로 이루어졌다. 매매 계약일 현재 동 오피스텔의 시장 순임대료(순영업소득)는 연간 3,000만원이고, 임대차계약 만료시 재매도 가치는 65,000만원으로 예상되고 있다. 제시된 자료를 활용하여 다음의 물음에 답하시오. (10점) \par
\medskip\noindent 물음 1) 이 매매사례를 평가에 채택할 경우 사정보정률(백분율로 표시하되, 소수점 이하에서 반올림)을 산출하시오. (5점) \par
\medskip\noindent 물음 2) 시장가치가 동일할 경우, 계약임대료(순영업소득)에 기한 환원율과 시장임대료(순영업소득)에 기한 환원율과의 차이를 산출하고 이 차이가 의미하는 바가 무엇인지 약술하시오. (5점) \par
\end{longquestion}
